\section{Analisis Kebutuhan dan Desain}

Kebutuhan diturunkan langsung dari peran nyata yang setiap hari memakai Atlas di Laboratorium MGM.

\subsection{Tiga peran pengguna utama}

\begin{description}
\item[Admin instansi.] Mengatur metode masuk, integrasi, dan kebijakan seluruh instans lewat Godmode. Butuh kendali penuh tanpa harus membaca kode.
\item[Ketua proyek / koordinator.] Membagi pekerjaan lewat papan Kanban atau linimasa, memantau progres, dan menjaga proyek tidak macet di satu orang.
\item[Anggota / kontributor.] Perlu tahu tugas apa yang jadi tanggung jawabnya hari ini, dan bisa berkoordinasi tanpa berpindah aplikasi.
\end{description}

\subsection{Kebutuhan fungsional dan non-fungsional}

Tabel berikut merangkum kebutuhan lintas enam kategori, dari fungsional hingga portabilitas deployment.

\begin{table*}[t]
\centering
\caption{Kebutuhan fungsional dan non-fungsional Atlas.}
\label{tab:kebutuhan}
\small
\begin{tabular}{@{}L{3.2cm}L{13cm}@{}}
\toprule
\textbf{Kategori} & \textbf{Kebutuhan} \\
\midrule
Fungsional & Manajemen tugas multi-tampilan (Kanban/List/Timeline), obrolan real-time, panggilan suara, catatan \& whiteboard kolaboratif, notifikasi terpadu, manajemen peran dan izin \\
Otentikasi & Minimal enam metode masuk termasuk SSO institusi (SAML/OIDC), tanpa biaya tambahan pada jumlah pengguna berapa pun \\
Kinerja & Sinkronisasi dokumen dan suara real-time tidak melewati backend sebagai bottleneck (koneksi langsung peer-to-sidecar) \\
Keamanan & Sesi tervalidasi di basis data per permintaan, kata sandi di-\textit{hash} dengan bcrypt, rahasia godmode dienkripsi AES-256-GCM \\
Aksesibilitas & Kontras warna lolos audit WCAG otomatis di 24 tema $\times$ 2 mode, target sentuh minimal 36px \\
Portabilitas & Dapat di-\textit{deploy} lewat Docker Compose di infrastruktur mana pun, tidak terkunci pada satu penyedia cloud \\
\bottomrule
\end{tabular}
\end{table*}

\FloatBarrier
\subsection{Desain sebagai kebutuhan, bukan pemanis}

Sistem desain MGM Laboratory (token warna tertutup, tipografi Bricolage Grotesque + Geist, motif geometris satu warna per permukaan) diterapkan konsisten dari landing page hingga panel admin Atlas. Ini bukan keputusan kosmetik: produk yang dipakai 65+ orang setiap hari perlu terasa dapat dipercaya dan konsisten, sama seperti alasan tim akhirnya berhenti memakai kombinasi alat yang tampilannya saling bertabrakan.
